
Across 4,493 open-weight language models submitted to the Open LLM Leaderboard, only 3.5\% document any pretraining data curation practice in their model cards. Naive OLS regression indicates that those models score lower on standardized benchmarks than their undocumented counterparts (beta = -3.45, p = 1.1 x $10^-8$). This counterintuitive result reflects a size-vintage confound rather than a causal effect, and it motivates the systematic registry introduced in this paper.

Data quality has been formalized as a determinant of language model loss. Subramanyam et al. (2025) extend the Chinchilla scaling law \citep{Hoffmann2022Training} to include a dimensionless data quality scalar Q, yielding a loss function L(N, D, Q) with gamma\_CLM approximately 0.39. However, directly measuring Q requires controlled noise injection experiments or access to proprietary training corpora. For the majority of deployed open-weight models, neither is available.

Model cards \citep{Mitchell2019Model} and datasheets \citep{Gebru2021Datasheets} are structured transparency documents that describe training data practices. When model developers document deduplication, domain composition choices, or decontamination procedures, they provide binary indicators of curation practice that are observable at scale without corpus access. Whether these documentation indicators serve as reliable proxies for data quality -- whether they predict benchmark performance beyond scale -- has not been empirically tested. No systematic registry linking model card curation documentation to benchmark performance at scale has previously existed.

This gap is relevant to AI governance. If documentation correlates with quality, transparency frameworks carry evidential weight for model evaluation. If the correlation is absent or confounded, governance frameworks that treat documentation as a quality proxy lack empirical support.

This work makes three contributions:

\begin{enumerate}
\item \textbf{The LLM Documentation-Benchmark Registry (n = 4,493).} The first large-scale dataset joining binary curation documentation indicators from HuggingFace model cards with Open LLM Leaderboard v2 benchmark scores across six evaluation tasks (IFEval, BBH, MATH Lvl 5, GPQA, MUSR, MMLU-PRO). Three of four documentation features (deduplication, domain composition, decontamination) exhibit non-zero variance.
\end{enumerate}

\begin{enumerate}
\item \textbf{Targeted family sampling methodology.} An initial alphabetical retrieval approach (H-E1) yielded only 177 model cards concentrated in the 0-A range, missing all major documented families. A targeted family sampling approach (H-E1-v2) prioritizes LLaMA, Mistral, Qwen, Falcon, Pythia, and OLMo families, retrieving 3,749 cards and achieving sufficient feature variance for regression analysis.
\end{enumerate}

\begin{enumerate}
\item \textbf{Confound identification.} Naive OLS produces a statistically significant negative documentation coefficient (beta = -3.45, p = 1.1 x $10^-8$). The most plausible explanation is a size-vintage confound: well-documented base models are disproportionately older and smaller-parameter models, while high-scoring V2 models tend to be newer, larger, or instruction-tuned derivatives. Additionally, perplexity filtering is absent from model cards under its canonical name, revealing a vocabulary gap in current documentation standards.
\end{enumerate}

The paper is organized as follows. Section 2 surveys related work. Section 3 describes the registry construction methodology. Section 4 presents the experimental design. Section 5 reports results. Section 6 discusses implications and limitations. Section 7 concludes.
